\chapter{Airbnb 算法心法}

这一章不讲具体题目，只讲“做 Airbnb 算法题时脑子里应该装着什么”。把这套心法装进去，再看下一章的十道题，你会发现它们其实是同一套思维的不同投影。

\section{Airbnb 算法题的三个特征}

\begin{enumerate}
  \item \textbf{偏模拟与建模，少冷门 DP}。Airbnb 爱出“把一段现实逻辑翻译成代码”的题：文本排版、倒水、分页去重。这类题没有花哨算法，但边界极多，考的是你能不能把规则\emph{不重不漏}地落到代码上。
  \item \textbf{多步迭代，会追加需求}。面试官常在你写完后说“现在再加一个条件……”。所以第一版就要写得\textbf{可扩展}：把规则抽成函数、把状态封进对象，而不是一坨写死的 \code{if/else}。
  \item \textbf{重视面向对象设计（OOD）}。迭代器、分页器、限流器这类题，本质是“设计一个类”，考的是接口划分与封装。
\end{enumerate}

\begin{idea}[为“被追加需求”而写代码]
区分“能 AC 的代码”和“Airbnb 想要的代码”的关键，是\textbf{可扩展性}。一个实用的习惯：把“业务规则”从“控制流程”里抽出来。例如分页去重题，不要把“同一 host 不相邻”这条规则硬编码在循环里，而是写成一个 \code{bool canPlace(item, page)} 判定函数。当面试官追加“VIP host 优先”时，你只需改一个函数，而不是重写整个循环。这一个习惯，往往就是 hire 和 strong hire 的差别。
\end{idea}

\section{解题的标准动作（边写边说）}

\begin{enumerate}
  \item \textbf{澄清}：先复述题意、确认输入范围与边界（空输入？重复元素？大小限制？），举一个自己的小样例。
  \item \textbf{先暴力，再优化}：大方说出暴力解和它的复杂度，再讲怎么优化。面试官想看到你的思考路径，而不是凭空蹦出最优解。
  \item \textbf{先讲思路，再写码}：达成一致后再动手，避免写到一半推倒重来。
  \item \textbf{边写边解释}：让面试官跟得上你的每一步。
  \item \textbf{自测}：写完用样例和边界 case 走一遍，主动发现 bug。
\end{enumerate}

\begin{keypoint}[复杂度报价]
写完一定要主动报时间/空间复杂度，并说明瓶颈在哪。这是基本职业素养，不报会扣分。
\end{keypoint}

\section{必须烂熟于心的工具箱}

下面这些是 Airbnb 高频题反复用到的“原子工具”。下一章每道题都会从这个工具箱里取件：

\begin{center}
\begin{tabularx}{\linewidth}{@{}l X@{}}
\toprule
\textbf{工具} & \textbf{典型用途（对应下一章的题）} \\
\midrule
哈希表 & 时间键值存储、回文对、火星词典建图 \\
对顶堆（双堆） & 数据流中位数 \\
二分查找 & 时间键值存储（在有序版本上找 $\le t$ 的最大值） \\
回溯 & 组合总和 \\
BFS / 状态编码 & 滑动谜题（状态空间搜索） \\
拓扑排序 & 火星词典（字符偏序） \\
栈 & 嵌套迭代器的惰性展开 \\
模拟 + 规则函数 & 文本对齐、倒水、分页去重 \\
\bottomrule
\end{tabularx}
\end{center}

\begin{followup}[“能不能让它支持流式/超大输入？”]
这是 Airbnb 最爱的追问方向之一。应答套路：把“一次性加载全部数据”的解法，改造成\textbf{迭代器/在线算法}——只保留必要状态。例如中位数题天然是流式（双堆增量维护）；时间键值存储可以谈“如果数据放不下内存，按 key 分片或落盘 + 二分”。提前想好这条线，追问来了就能从容接住。
\end{followup}

\begin{pitfall}[只顾最优解，丢了沟通]
Airbnb 算法轮里，一个全程沉默、最后甩出最优解的候选人，评分往往不如一个边讲边写、思路清晰但代码略有瑕疵的候选人。\textbf{沟通本身就是评分项}。
\end{pitfall}
